\scanpage{370}
\mspage{109}
% archive sheet number: 369
une $\mathrm{Hot}$-équivalence. Sinon, cela prouvera que dans le cas général, le \deleted{type} $\mathrm{Hot}$-type de $C$ (et a fortiori \deleted{\unsure{la}}\unsure{sa|la} classe d'isom. dans $\mathrm{Cat\,hot}$) n'a pas de sens géométrique intéressant. \deleted{\unsure{C}} C'est à dire que $C$ est intéressant uniquement en \uline{\unsure{tant} que catégorie fibrée sur $T$}, permettant \inserted{notamment} de construire l'\inserted{\unsure{ind-}}objet $H_T$ de $\mathrm{Cat\,hot}$, ou de $\mathrm{Hot}$, mais pas en tant que ``catégorie nue''.

Mais je viens de me convaincre que $C_{\mathbb N^\circ}\mathrel{\unsure{-|\to}}C_\infty$ est bien une $\mathrm{Hot}$-équivalence. Donc la flèche $C_{\mathbb N^\circ}\longrightarrow C$ définit une flèche \uline{dans $\mathrm{Hot}$}
\[
C_\infty\longrightarrow C
\]
(en composant avec l'inverse \inserted{dans $\mathrm{Hot}$} de la précédente). Est-ce un iso ? Donc pour cette question, on est ramené à la question pour $C_{\mathbb N^\circ}\longrightarrow C$. Mais si on prend $C=T$, cela pose déjà la même question pour $\underline{\mathbb N}^\circ\longrightarrow\deleted{\unsure{C_\infty}}T$, donc la question si $T$ est asphérique, sous la seule hypothèse que $T_0$ soit filtrante à g. et dénombrable à l'infini \inserted{(à g.)}. Mais bien sûr c'est faux (p. ex. si $T$ n'a qu'un seul objet, il n'est pas pour autant asphérique !).

Cela me convainc donc qu'à défaut
